\section{Annotated-Relation Stress Test: Pedestrian Groups}
\label{sec:atr-transfer}

\paragraph{Role in the evidence chain.}
The Chicago and NYC releases contain real rows but omit the operational
relation, so they cannot measure whether a feasible-world frontier contains the
realized world. Controlled data provide exact truth but may encode the method's
own assumptions. ATR occupies the narrow middle layer: we apply EventFrontier
to DIAMOR
pedestrian data, which provide timestamped positions and motion attributes from
public corridors together with manually coded social groups
\citep{atrgroupdata2015,zanlungo2014potential}.  The tracking rows are produced
by range-sensor systems, whereas group membership is coded separately.  This
creates an annotated relation that can be masked from candidate edge scoring
and retained for a controlled audit. The purpose is not to claim general
external validity. It is to stress-test the candidate-support gate on real row
relations and measure what fails when that gate omits annotated edges.

\paragraph{Release operator and scope.}
We retain timestamp, planar position, speed, and direction of motion.  At each
fixed snapshot, two visible rows are candidate partners when their separation
is at most $r\in\{1,1.5,2,3,5\}$ meters and their directions differ by at most
$90$ degrees.  The true number $q$ of visible dyads is disclosed, but their
memberships are hidden.  Labels also determine eligibility and which annotated
larger or uncertain groups are excluded.  Candidate edge scoring itself is
label-free, but the full experiment is therefore an oracle-assisted,
$q$-conditioned masked-label audit.  We conservatively quarantine recoverable
positive IDs in partial, nonpositive, multiple-ID, one-sided, or conflicting
annotations rather than silently treat them as singletons.  Each remaining
feasible world is a cardinality-$q$ matching, and the query is its mean
within-dyad separation.  The deterministic point comparator selects a
minimum-separation matching.

The v4 evidence protocol fixes 136 snapshots on a two-minute grid for each day.
DIAMOR-1 is the pilot day.  DIAMOR-2 applies the same grid and candidate rules as
a predeclared follow-up; because it had already been touched by an earlier
sparse audit, we do not describe it as a pristine holdout.

\begin{table}[t]
\caption{ATR dyad audit. Coverage is evaluated only where every annotated dyad
edge is present in candidate support. ``Flagged'' counts point-decision errors
occurring at thresholds crossed by the frontier.}
\label{tab:atr-truth}
\centering
\small
\begin{tabular}{@{}lrrrr@{}}
\toprule
Day & Eligible snapshots & Cells & Truth coverage & Errors flagged \\
\midrule
DIAMOR-1 (pilot) & 88 & 440 & 388/388 & 36/36 \\
DIAMOR-2 (follow-up) & 82 & 410 & 381/381 & 32/32 \\
\midrule
Total & 170 & 850 & 769/769 & 68/68 \\
\bottomrule
\end{tabular}
\end{table}

\paragraph{Truth coverage and candidate completeness.}
All 850 radius--snapshot cells terminate with either replayed numerical MILP
optima (806 cells) or candidate-graph infeasibility (44 cells).  All 195 cells
within a 12-vertex limit also agree with an independent subset-recursion oracle.
Replay checks the raw solution's integrality, degree, cardinality, objective,
and primal--dual residuals before an endpoint is published.
Conditional on the annotated matching being representable, the frontier covers
the realized mean separation in all $769/769$ cells.  The closest-pair point
rule makes 68 threshold errors in those cells; every error occurs where the
frontier crosses the threshold. These conditional coverage and flagging counts
are positive controls implied by optimizing over a feasible set that contains
the annotated world. They check implementation, but are not an empirical
success rate or independent evidence of transfer.

Candidate completeness is not automatic.  Across both days, a 1-meter radius
contains the annotated world in only 124 of 170 eligible snapshots; this rises
to 157 at 1.5 meters and 165 at 5 meters.  The gain comes at the cost of wider
identified sets: on the predeclared follow-up day, mean width rises from
$0.043$ meters at $r=1$ to $3.861$ meters at $r=5$.  At $r=5$, 410 of 495
valid-support threshold decisions (82.8\%) remain ambiguous.  Scalar aggregate
coverage under a misspecified graph is not counted as relational coverage: the
true value can fall inside an interval by coincidence even when the annotated
matching itself is absent.  Under misspecified support, 16 wrong point
decisions are not flagged, across eight snapshots.

\paragraph{Removing oracle cardinality.}
We next union the fixed-$q$ frontiers over either $q\pm1$ or every positive
cardinality through half the visible population. The latter removes labels
from endpoint cardinality, although labels still determine eligibility,
exclusions, and evaluation. Relative to oracle $q$, the all-positive-$q$
regime makes 826 rather than 806 cells feasible because some candidate graphs
admit a smaller matching than the annotated count. This is apparent
feasibility, not restored relational coverage. Mean width increases from
$1.176$ to $1.297$ meters, and certified threshold decisions fall from
$1286/2418$ (53.2\%) to $1136/2478$ (45.8\%). All 20 false certificates occur
where candidate support omits the annotated world; none occurs under
representable truth.

\paragraph{Cross-day, multi-query point reconstruction.}
To avoid aligning every point rule with its evaluated query, we train a
class-balanced logistic dyad score on DIAMOR-1 using only distance, heading
difference, and speed difference, freeze it, and construct cardinality-$q$
matchings on DIAMOR-2. The same matching is evaluated for mean distance, mean
speed gap, and mean heading gap. Across the 381 exact cells with representable
truth, the learned and closest-distance rules have similar mean dyad recall
(89.24\% versus 89.62\%) but different downstream errors. The learned rule
makes respectively 44, 70, and 35 threshold errors for the three queries,
compared with 32, 44, and 38 for closest-distance; every error is
frontier-ambiguous. Thus comparable edge recovery does not imply comparable
decision reliability. The linear scorer is a transparent cross-day baseline,
not a state-of-the-art trajectory model.

\paragraph{Calibrating candidate support across days.}
Finally, on DIAMOR-1 we search a fixed $8\times6$ radius--motion-angle grid,
require 95\% snapshot-level relational coverage, and select the qualifying
rule with the fewest mean candidate edges. The selected 3-meter, 120-degree
rule contains the annotated world in 85/88 pilot snapshots with 4.52 candidate
edges per snapshot. Frozen on DIAMOR-2, it covers 81/82 snapshots, but candidate
density rises to 12.17 edges per snapshot and only 29/246 threshold decisions
are certified. Thus coverage transfers in this comparison while sparsity and
decision informativeness do not. This is not a universal support rule: the
density shift motivates reporting both relational coverage and downstream
certification yield.

A post-hoc mechanism audit then caps each node at its nearest incident edges
inside that frozen support. DIAMOR-1 selects $k=2$ as the sparsest cap preserving
the fixed rule's 85/88 pilot coverage. On DIAMOR-2 it preserves 81/82 coverage
and removes 80/998 candidate edges (8.0\%), yet certifies exactly the same
29/246 decisions. Simple degree control therefore reduces graph bulk without
restoring informativeness; structural alternative matchings survive the cap.
An endpoint-witness decomposition makes this failure more concrete. All 82
fixed graphs have positive width, and 73 returned lower--upper witness pairs
differ through multiple disconnected alternating components. These components
are small---at most four edges---and overwhelmingly paths. The cap leaves every
lower endpoint unchanged and contracts 21 upper endpoints (0.177 meters on
average among changed cells), but crosses no evaluated threshold. Graph
compression and decision identification therefore separate in this audit.
This is the empirical pattern allowed by
Proposition~\ref{prop:k2thresholdcontraction}: a nested candidate graph can
strictly contract continuous endpoints while leaving every evaluated threshold
inside the contracted frontier. The small-component counts illustrate
Theorem~\ref{thm:k2alternating}; they are not required for its conclusion.

\paragraph{What ATR adds---and what it does not.}
The informative result is the failure boundary, not $769/769$ conditional
coverage. Real annotated dyads show that a plausible geometric support rule can
omit the joint world even when optimization remains feasible; in those
misspecified cells, frontier ambiguity need not protect the analyst. The radius
sweep also exposes a coverage--informativeness trade-off: wider support retains
more annotated worlds while rapidly making threshold decisions ambiguous.
Synthetic experiments isolate this mechanism under known truth; ATR shows that
the same support failure is present in a separately collected relation
annotation. Chicago and NYC then assess release-specific scale and ambiguity,
not truth coverage.

This audit validates conditional coverage for fixed snapshots, dyads, a
disclosed support count, label-assisted cohort construction, and one geometric
query.  The group labels were produced by one coder and are treated as an
annotated reference rather than error-free physical truth.  It neither
estimates the population prevalence of social groups nor validates the Chicago
or NYC candidate rules.  ATR data are research-use-only and are not
redistributed.  Full witnesses remain in a controlled local audit; the public
artifact stores the protocol, hashes, aggregate summaries, and statuses without
row IDs or relation assignments.
